\section{第 2.1 节教材习题与订正}
\label{exercise-set:QM-C02-S01}

\exercisemeta
  {QM-C02-S01-EX}
  {2026-08-27}
  {Shankar，第 2 章 \S 2.1，练习 2.1.1--2.1.3}
  {完成三道教材习题并逐题核对}
  {题意、Euler--Lagrange 推导与关键计算已核对}

\exercisequestion{QM-C02-S01-E01}{教材练习 2.1.1：一维简谐振子}

\textbf{对应知识点：}\relatedtopic{QM-C02-S01-T02}{Euler--Lagrange 方程的完整变分推导}。

\textbf{题意。}质量为 $m$ 的物块位于无摩擦水平面上，并与劲度系数为 $k$ 的弹簧相连。写出 Lagrangian 并求运动方程。教材位置：印刷页 p.~83，PDF p.~96。

\begin{checkedsolution}
取物块偏离平衡位置的位移为 $x$。动能、势能和 Lagrangian 分别为
\[
  T=\frac12m\dot x^2,
  \qquad
  V=\frac12kx^2,
  \qquad
  \mathcal L=\frac12m\dot x^2-\frac12kx^2.
\]
所需偏导为
\[
  \frac{\partial\mathcal L}{\partial\dot x}=m\dot x,
  \qquad
  \frac{\partial\mathcal L}{\partial x}=-kx.
\]
代入 Euler--Lagrange 方程得到
\begin{equation}
  \boxed{m\ddot x+kx=0,}
  \qquad
  \omega=\sqrt{\frac{k}{m}}.
  \label{eq:QM-C02-S01-E01-oscillator}
\end{equation}
这与 Newton 方程 $m\ddot x=-kx$ 完全相同。
\end{checkedsolution}

\textbf{检查点：}势能进入 $\mathcal L=T-V$ 时带负号，因此 $\partial\mathcal L/\partial x=-kx$；最终方程的恢复力方向必须与位移相反。

\exercisequestion{QM-C02-S01-E02}{教材练习 2.1.2：耦合质量块的 Lagrangian}

\textbf{对应知识点：}\relatedtopic{QM-C02-S01-T03}{广义坐标、正则动量与循环坐标}，以及第 1.8 节的耦合振子简正模。

\textbf{题意。}两个相同质量 $m$ 的物块沿直线运动，三个相同弹簧 $k$ 依次连接墙、物块 1、物块 2 与另一侧墙。令 $x_1,x_2$ 为两个物块相对平衡位置的同向位移。写出 Lagrangian，用 Euler--Lagrange 方程推导耦合运动方程，并与第 1.8 节的 Newton 方程比较。教材位置：印刷页 p.~83，PDF p.~96。

\begin{checkedsolution}
两个物块的总动能为
\begin{equation*}
  T=\frac12m(\dot x_1^2+\dot x_2^2).
\end{equation*}
三个弹簧的有向形变量可取为 $x_1$、$x_2-x_1$ 与 $-x_2$；势能只依赖平方，因此
\begin{align*}
  V
  &=\frac12kx_1^2
    +\frac12k(x_2-x_1)^2
    +\frac12kx_2^2\\
  &=k(x_1^2+x_2^2-x_1x_2).
\end{align*}
所以
\begin{equation}
  \mathcal L
  =\frac12m(\dot x_1^2+\dot x_2^2)
  -\frac12k\left[x_1^2+(x_2-x_1)^2+x_2^2\right].
  \label{eq:QM-C02-S01-E02-lagrangian}
\end{equation}
对 $x_1$，
\[
  \frac{\partial\mathcal L}{\partial\dot x_1}=m\dot x_1,
  \qquad
  \frac{\partial\mathcal L}{\partial x_1}=-(2kx_1-kx_2),
\]
故
\[
  m\ddot x_1+2kx_1-kx_2=0.
\]
同理，对 $x_2$ 有
\[
  \frac{\partial\mathcal L}{\partial\dot x_2}=m\dot x_2,
  \qquad
  \frac{\partial\mathcal L}{\partial x_2}=-(2kx_2-kx_1),
\]
从而
\[
  m\ddot x_2+2kx_2-kx_1=0.
\]
矩阵形式为
\begin{equation}
  \boxed{
  m\begin{pmatrix}\ddot x_1\\\ddot x_2\end{pmatrix}
  +k\begin{pmatrix}2&-1\\-1&2\end{pmatrix}
  \begin{pmatrix}x_1\\x_2\end{pmatrix}=0.}
  \label{eq:QM-C02-S01-E02-coupled-equations}
\end{equation}
这正是第 1.8 节由逐个分析弹簧力得到的方程。作为独立检查，刚度矩阵在 $(1,1)^T$ 和 $(1,-1)^T$ 上的特征值分别为 $k$ 和 $3k$，故
\[
  \omega_I=\sqrt{\frac{k}{m}},
  \qquad
  \omega_{II}=\sqrt{\frac{3k}{m}},
\]
与原简正模结果一致。
\end{checkedsolution}

\textbf{检查点：}中间弹簧的形变量必须是相对位移 $x_2-x_1$。若把它误写成 $x_1+x_2$，同向等位移模式会错误地拉伸中间弹簧。

\exercisequestion{QM-C02-S01-E03}{教材练习 2.1.3：三维中心势的球坐标方程}

\textbf{对应知识点：}\relatedtopic{QM-C02-S01-T05}{三维中心势中的球坐标与角动量}。

\textbf{题意。}质量为 $m$ 的粒子在三维中心势 $V(r)$ 中运动。使用球坐标写出 Lagrangian，求 $r,\theta,\phi$ 的 Euler--Lagrange 方程，并判断循环坐标与守恒量。教材位置：印刷页 p.~83，PDF p.~96。

\begin{checkedsolution}
球坐标中的速度平方为
\[
  v^2=\dot r^2+r^2\dot\theta^2
  +r^2\sin^2\theta\,\dot\phi^2.
\]
因此
\begin{equation}
  \mathcal L
  =\frac12m\left(
    \dot r^2+r^2\dot\theta^2
    +r^2\sin^2\theta\,\dot\phi^2
  \right)-V(r).
  \label{eq:QM-C02-S01-E03-lagrangian}
\end{equation}

对 $r$，
\[
  \frac{\partial\mathcal L}{\partial\dot r}=m\dot r,
  \qquad
  \frac{\partial\mathcal L}{\partial r}
  =mr\dot\theta^2+mr\sin^2\theta\,\dot\phi^2-V'(r),
\]
所以
\begin{equation}
  \boxed{
  m\ddot r-mr\dot\theta^2
  -mr\sin^2\theta\,\dot\phi^2+V'(r)=0.}
  \label{eq:QM-C02-S01-E03-radial}
\end{equation}

对 $\theta$，
\[
  \frac{\partial\mathcal L}{\partial\dot\theta}=mr^2\dot\theta,
  \qquad
  \frac{\partial\mathcal L}{\partial\theta}
  =mr^2\sin\theta\cos\theta\,\dot\phi^2,
\]
所以
\begin{equation}
  \boxed{
  r^2\ddot\theta+2r\dot r\dot\theta
  -r^2\sin\theta\cos\theta\,\dot\phi^2=0.}
  \label{eq:QM-C02-S01-E03-theta}
\end{equation}

$\phi$ 不显含于完整的 $\mathcal L$，因此
\begin{equation}
  \boxed{
  \frac{d}{dt}\left(
    mr^2\sin^2\theta\,\dot\phi
  \right)=0,}
  \label{eq:QM-C02-S01-E03-phi}
\end{equation}
即
\begin{equation}
  \boxed{
  p_\phi=mr^2\sin^2\theta\,\dot\phi=L_z
  =\text{constant}.}
  \label{eq:QM-C02-S01-E03-lz}
\end{equation}
$r$ 显含于动能和势能，$\theta$ 通过 $\sin^2\theta$ 显含于动能，所以二者都不是循环坐标。中心势虽只直接让 $\phi$ 成为这一坐标系中的循环坐标，但完整旋转对称性仍保证整个角动量向量守恒；当前球坐标只是让 $L_z$ 的守恒最明显。
\end{checkedsolution}

\textbf{检查点：}循环坐标的判据是 $\partial\mathcal L/\partial q_i=0$，不是 $\partial V/\partial q_i=0$。$V(r)$ 与 $\theta$ 无关仍不足以让 $\theta$ 成为循环坐标，因为球坐标动能含 $\sin^2\theta$。

\section{第 2.2 节自编检验题与订正}
\label{exercise-set:QM-C02-S02}

\exercisemeta
  {QM-C02-S02-EX}
  {2026-08-28}
  {Shankar，第 2 章 \S 2.2；教材本节无正式例题或习题}
  {完成两道 self-contained check examples（自编检验题）}
  {题目来源类型、单位制、推导与物理解释已核对}

\exercisequestion{QM-C02-S02-E01}{自编检验题 1：均匀磁场中的正则动量}

\textbf{对应知识点：}\relatedtopic{QM-C02-S02-T03}{从 Lagrangian 完整恢复 Lorentz force}与\relatedtopic{QM-C02-S02-T04}{正则动量、机械动量与广义势}。

\textbf{题意。}在 Gaussian units 中，带电量 $q$、质量 $m$ 的粒子在 $xy$ 平面运动。取 $\phi=0$ 与 Landau gauge
\[
  \bm A=Bx\,\hat{\bm y},
\]
它对应均匀磁场 $\bm B=B\hat{\bm z}$。写出 Lagrangian，求正则动量与运动方程，并判断哪个动量分量守恒。

\begin{checkedsolution}
因为 $\bm v\cdot\bm A=Bx\dot y$，所以
\begin{equation}
  \mathcal L
  =\frac12m(\dot x^2+\dot y^2)+\frac{qB}{c}x\dot y.
  \label{eq:QM-C02-S02-E01-lagrangian}
\end{equation}
正则动量为
\begin{equation}
  p_x=m\dot x,
  \qquad
  p_y=m\dot y+\frac{qB}{c}x.
  \label{eq:QM-C02-S02-E01-canonical-momenta}
\end{equation}
对 $x$ 使用 Euler--Lagrange 方程：
\[
  \frac{d}{dt}(m\dot x)-\frac{qB}{c}\dot y=0,
\]
故
\[
  m\ddot x=\frac{qB}{c}\dot y.
\]
$y$ 不显含于 $\mathcal L$，因此是循环坐标：
\[
  \frac{d}{dt}\left(m\dot y+\frac{qB}{c}x\right)=0,
\]
即
\[
  m\ddot y=-\frac{qB}{c}\dot x.
\]
合并为
\begin{equation}
  m\ddot{\bm r}
  =\frac qc\dot{\bm r}\times B\hat{\bm z},
  \label{eq:QM-C02-S02-E01-lorentz}
\end{equation}
与 Lorentz force 一致。守恒的是 canonical momentum $p_y$，而 mechanical momentum 的 $y$ 分量 $m\dot y$ 一般不守恒；两者相差 $qBx/c$。
\end{checkedsolution}

\textbf{检查点：}循环坐标直接保证其共轭正则动量守恒，而不是自动保证相应的 Cartesian mechanical momentum 分量守恒。

\exercisequestion{QM-C02-S02-E02}{自编检验题 2：固定长度绳子的约束力}

\textbf{对应知识点：}\relatedtopic{QM-C02-S02-T05}{相互作用、坐标几何、约束与非保守力}。

\textbf{题意。}质量为 $m$ 的小球在平面内运动，并由固定在原点、长度为 $R$ 的理想不可伸长轻绳约束。系统没有普通势能。分别用消去约束坐标与 Lagrange multiplier 两种方法写出方程，并说明为什么 $V=0$ 不等于张力不存在。

\begin{checkedsolution}
若直接使用约束 $\rho=R$，唯一独立坐标为 $\phi$，约化 Lagrangian 是
\begin{equation}
  \mathcal L_{\mathrm{red}}=\frac12mR^2\dot\phi^2.
  \label{eq:QM-C02-S02-E02-reduced-lagrangian}
\end{equation}
因此
\[
  \frac{d}{dt}(mR^2\dot\phi)=0,
\]
角动量守恒。张力沿径向，而允许虚位移沿切向，所以理想张力不对允许虚位移做功；消去 $\rho$ 后，它不出现在 $\phi$ 的方程中。

若还要计算张力，则保留 $\rho$ 并引入乘子 $\lambda$：
\begin{equation}
  \mathcal L'
  =\frac12m(\dot\rho^2+\rho^2\dot\phi^2)
  -\lambda(\rho-R).
  \label{eq:QM-C02-S02-E02-multiplier-lagrangian}
\end{equation}
径向方程为
\[
  m\ddot\rho-m\rho\dot\phi^2+\lambda=0.
\]
代入 $\rho=R$ 与 $\ddot\rho=0$，得到
\begin{equation}
  \boxed{\lambda=mR\dot\phi^2.}
  \label{eq:QM-C02-S02-E02-tension}
\end{equation}
$\lambda$ 是绳子张力的大小。$V=0$ 只表示没有写入普通势能的相互作用，不能删除独立给定的约束及其反作用力。
\end{checkedsolution}

\textbf{检查点：}若约束坐标已被消去，约束力通常不会出现在剩余 Euler--Lagrange 方程中；要恢复约束力，需要保留约束并使用乘子，或另作受力分析。

\section{第 2.3 节教材习题与订正}
\label{exercise-set:QM-C02-S03}

\exercisemeta
  {QM-C02-S03-EX}
  {2026-08-29}
  {Shankar，第 2 章 \S 2.3，练习 2.3.1}
  {完成一道教材习题并核对坐标代换}
  {速度变换、交叉项抵消、约化质量系数与最终 Lagrangian 已核对}

\exercisequestion{QM-C02-S03-E01}{教材练习 2.3.1：分离二体 Lagrangian}

\textbf{对应知识点：}\relatedtopic{QM-C02-S03-T02}{动能分离与约化质量}。

\textbf{题意。}从两个相互作用粒子的原 Lagrangian 出发，改用质心坐标 $\bm R$ 与相对坐标 $\bm r$，推导分离后的 Lagrangian。教材位置：印刷页 p.~86，PDF p.~99。

\begin{checkedsolution}
原式为
\begin{equation}
  \mathcal L
  =\frac12m_1|\dot{\bm r}_1|^2
  +\frac12m_2|\dot{\bm r}_2|^2
  -V(\bm r_1-\bm r_2).
  \label{eq:QM-C02-S03-E01-original}
\end{equation}
令
\[
  M=m_1+m_2,
  \qquad
  \bm r=\bm r_1-\bm r_2,
  \qquad
  \bm R=\frac{m_1\bm r_1+m_2\bm r_2}{M}.
\]
逆变换及其时间导数为
\begin{align*}
  \bm r_1&=\bm R+\frac{m_2}{M}\bm r,
  &\dot{\bm r}_1&=\dot{\bm R}+\frac{m_2}{M}\dot{\bm r},\\
  \bm r_2&=\bm R-\frac{m_1}{M}\bm r,
  &\dot{\bm r}_2&=\dot{\bm R}-\frac{m_1}{M}\dot{\bm r}.
\end{align*}
因此动能为
\begin{align*}
  T={}&\frac12m_1
  \left|
    \dot{\bm R}+\frac{m_2}{M}\dot{\bm r}
  \right|^2
  +\frac12m_2
  \left|
    \dot{\bm R}-\frac{m_1}{M}\dot{\bm r}
  \right|^2\\
  ={}&\frac12(m_1+m_2)|\dot{\bm R}|^2
  +\left(
    \frac{m_1m_2}{M}-\frac{m_1m_2}{M}
  \right)\dot{\bm R}\cdot\dot{\bm r}\\
  &+\frac12\left(
    \frac{m_1m_2^2}{M^2}
    +\frac{m_2m_1^2}{M^2}
  \right)|\dot{\bm r}|^2.
\end{align*}
中间的交叉项严格为零，而最后一个系数化简为
\begin{align*}
  \frac{m_1m_2^2+m_2m_1^2}{M^2}
  &=\frac{m_1m_2(m_1+m_2)}{M^2}
   =\frac{m_1m_2}{m_1+m_2}
   \equiv\mu.
\end{align*}
同时 $V(\bm r_1-\bm r_2)=V(\bm r)$，故
\begin{equation}
  \boxed{
  \mathcal L
  =\frac12M|\dot{\bm R}|^2
  +\frac12\mu|\dot{\bm r}|^2
  -V(\bm r),
  \qquad
  \mu=\frac{m_1m_2}{m_1+m_2}.}
  \label{eq:QM-C02-S03-E01-separated}
\end{equation}
\end{checkedsolution}

\textbf{检查点：}交叉项的抵消不是偶然删项，而是来自质心坐标中的质量权重。最终式中质心项的质量必须是 $m_1+m_2$，不能误写成 $m_1+m_1$；相对项的系数则是约化质量而非总质量。

\section{第 2.4 节自编检验题与订正}
\label{exercise-set:QM-C02-S04}

\exercisemeta
  {QM-C02-S04-EX}
  {2026-08-30}
  {Shankar，第 2 章 \S 2.4；教材本节无正式例题或习题}
  {完成一道 self-contained check example（自编检验题）}
  {平稳条件、局部运动规律、路径振幅与经典极限已核对}

\exercisequestion{QM-C02-S04-E01}{自编检验题 1：从平稳作用量到经典路径}

\textbf{对应知识点：}
\begin{itemize}[leftmargin=*]
  \item \relatedtopic{QM-C02-S04-T01}{全局变分不等于粒子预知未来}；
  \item \relatedtopic{QM-C02-S04-T02}{Stationary 不等于 minimum}；
  \item \relatedtopic{QM-C02-S04-T03}{路径积分预告与经典路径的显现}。
\end{itemize}

\textbf{题意。}回答下列概念问题：
\begin{enumerate}[label=(\arabic*),leftmargin=*]
  \item $\delta S=0$ 是否表示经典路径的作用量一定小于所有邻近路径？
  \item 固定未来终点是否表示粒子必须预知未来？
  \item 路径积分中各路径的复振幅模长相同，为什么经典极限仍突出平稳作用量路径？
\end{enumerate}

\begin{checkedsolution}
\begin{enumerate}[label=(\arabic*),leftmargin=*]
  \item 不是。$\delta S=0$ 只表示一阶变分为零，作用量可能在该路径处取得局部最小、局部最大或鞍点；还需检查更高阶变分才能进一步分类。
  \item 不表示。未来终点是边值问题中由求解者施加的条件。变分条件导出的 Euler--Lagrange 方程是局部微分方程，也可以结合当前位置和当前速度作为初值向未来演化，无需粒子获得未来信息。
  \item 相同的是普通实作用量下 $e^{iS[x]/\hbar}$ 的模长，不是相位或概率。非平稳路径附近的微小路径改变会造成 $O(\eta)$ 的作用量变化；当 $S/\hbar\gg1$ 时，相位迅速变化，邻近贡献大多相消。经典路径满足 $\delta S=0$，其附近的作用量差从 $O(\eta^2)$ 开始，相位变化较慢，因而一簇邻近路径的振幅能够更稳定地相干叠加。
\end{enumerate}
\end{checkedsolution}

\textbf{检查点：}路径积分中的“所有路径”讨论的是复振幅求和。经典路径的出现依赖 stationary phase，而不是量子理论先把其他路径删除。
